%% Generated by tools/gen_error_codes.py from docs/errors/ of the compiler;
%% edit the compiler's pages or tools/error_themes.txt, then rerun it.

\clearpage
\section{Mutability, aliasing and moves}
\label{sec:codes:memory}

The book's chapter \textit{Variables and memory management} specifies the rules behind these codes.

\begin{longtable}{@{}l p{0.78\linewidth} r@{}}
\toprule Code & Message & Page\\ \midrule \endhead
\hyperref[sec:codes:e4002]{E4002} & class alias operator :[ is prohibited on the left operand of an affectation & \pageref{sec:codes:e4002}\\
\hyperref[sec:codes:e4014]{E4014} & class alias operator :. is not usable to access fields & \pageref{sec:codes:e4014}\\
\hyperref[sec:codes:e4015]{E4015} & alias operator :. is only usable on class, record, entity, map or generator types, not \errhole{} & \pageref{sec:codes:e4015}\\
\hyperref[sec:codes:e4016]{E4016} & class alias operator :[ is only usable on class, record or entity types, not \errhole{} & \pageref{sec:codes:e4016}\\
\hyperref[sec:codes:e4029]{E4029} & movable type \errhole{} cannot be embedded within another type & \pageref{sec:codes:e4029}\\
\hyperref[sec:codes:e4045]{E4045} & discarding the constant property is prohibited & \pageref{sec:codes:e4045}\\
\hyperref[sec:codes:e4082]{E4082} & left operand of type \errhole{} is immutable & \pageref{sec:codes:e4082}\\
\hyperref[sec:codes:e4083]{E4083} & type must be immutable & \pageref{sec:codes:e4083}\\
\hyperref[sec:codes:e4088]{E4088} & implicit move of type \errhole{} is prohibited & \pageref{sec:codes:e4088}\\
\hyperref[sec:codes:e4089]{E4089} & ref of type \errhole{} must be mutable, but is const & \pageref{sec:codes:e4089}\\
\hyperref[sec:codes:e4092]{E4092} & cannot pass value of type \errhole{} as ref & \pageref{sec:codes:e4092}\\
\hyperref[sec:codes:e4099]{E4099} & index assignment of a value of type \errhole{} requires a mutable value & \pageref{sec:codes:e4099}\\
\hyperref[sec:codes:e4125]{E4125} & cannot alias an expand value, maybe alias and expand keywords are inverted? & \pageref{sec:codes:e4125}\\
\hyperref[sec:codes:e4126]{E4126} & cannot alias a lazy value, maybe alias and lazy keywords are inverted? & \pageref{sec:codes:e4126}\\
\hyperref[sec:codes:e4127]{E4127} & cannot copy an expand value, maybe copy and expand keywords are inverted? & \pageref{sec:codes:e4127}\\
\hyperref[sec:codes:e4128]{E4128} & cannot copy a lazy value, maybe copy and lazy keywords are inverted? & \pageref{sec:codes:e4128}\\
\hyperref[sec:codes:e4129]{E4129} & cannot copy an expand value, maybe dcopy and expand keywords are inverted? & \pageref{sec:codes:e4129}\\
\hyperref[sec:codes:e4130]{E4130} & cannot copy a lazy value, maybe dcopy and lazy keywords are inverted? & \pageref{sec:codes:e4130}\\
\hyperref[sec:codes:e4137]{E4137} & an iterator cannot be mutable, if it is not a reference or does not borrow mutable data & \pageref{sec:codes:e4137}\\
\hyperref[sec:codes:e4138]{E4138} & a parameter cannot be mutable, if it is not a reference or does not borrow mutable data & \pageref{sec:codes:e4138}\\
\hyperref[sec:codes:e4140]{E4140} & a lazy variable cannot be mutable, if it does not borrow mutable data & \pageref{sec:codes:e4140}\\
\hyperref[sec:codes:e4146]{E4146} & \errhole{} is not an aliasable type & \pageref{sec:codes:e4146}\\
\hyperref[sec:codes:e4151]{E4151} & not a lvalue & \pageref{sec:codes:e4151}\\
\hyperref[sec:codes:e4153]{E4153} & \errhole{} is a value iterator, and cannot be used as a lvalue & \pageref{sec:codes:e4153}\\
\hyperref[sec:codes:e4154]{E4154} & \errhole{} is a lazy variable, and cannot be used as a lvalue & \pageref{sec:codes:e4154}\\
\hyperref[sec:codes:e4155]{E4155} & \errhole{} is a map access, and cannot be used as a lvalue & \pageref{sec:codes:e4155}\\
\hyperref[sec:codes:e4156]{E4156} & \errhole{} is a value parameter, and cannot be used as a lvalue & \pageref{sec:codes:e4156}\\
\hyperref[sec:codes:e4157]{E4157} & value of type \errhole{} is not a lvalue & \pageref{sec:codes:e4157}\\
\hyperref[sec:codes:e4165]{E4165} & type \errhole{} is not a movable type & \pageref{sec:codes:e4165}\\
\hyperref[sec:codes:e4167]{E4167} & no copy exists for type \errhole{} & \pageref{sec:codes:e4167}\\
\hyperref[sec:codes:e4247]{E4247} & aliasing the value of type \errhole{} to create a constant borrowing is prohibited & \pageref{sec:codes:e4247}\\
\hyperref[sec:codes:e4248]{E4248} & aliasing the value of type \errhole{} to call begin and end iterator constant methods is useless & \pageref{sec:codes:e4248}\\
\hyperref[sec:codes:e4251]{E4251} & the move has no effect & \pageref{sec:codes:e4251}\\
\hyperref[sec:codes:e4252]{E4252} & referencing the value has no effect & \pageref{sec:codes:e4252}\\
\bottomrule
\end{longtable}

\errorcode{E4002}{class alias operator :[ is prohibited on the left operand of an affectation}
\label{sec:codes:e4002}

Raised during \textbf{validation}, from \texttt{ValidateErrorMessage::\allowbreak{}ALIAS\_\allowbreak{}INDEX\_\allowbreak{}AFFECT} (\textit{src/\allowbreak{}ymirc/\allowbreak{}semantic/\allowbreak{}validator/\allowbreak{}errors.yr}).

\subsubsection*{What it means}

The class alias index operator \tokennolst{:[} was used on the left of an assignment. That operator produces a borrowed view of the indexed value; the left of an assignment is storage, not a value, so there is nothing for the alias to borrow.

\subsubsection*{Example}

\textit{test\_\allowbreak{}resources/\allowbreak{}lit\_\allowbreak{}class/\allowbreak{}operators/\allowbreak{}test16.yr}:

%% check: skip
\begin{lstlisting}[style=coloredverbatimError]
class A {

    pub self () {}

    pub fn opIndexAssign (self, _ : i32, _ : i32) {}

}

fn main () {
    let dmut a = copy A ();
    a:[0] = 0;    
    a:.opIndexAssign (0, 0);
}
\end{lstlisting}

\begin{lstlisting}[style=bashVerb, breaklines=true, escapechar=~]
Error[E4002] : class alias operator :[ is prohibited on the left operand of an affectation
 --> test_resources/lit_class/operators/test16.yr:(11,6)
11  ~\lstbox{\textSFxi{}}~     a:[0] = 0;
    ~\lstbox{\textSFv{}}~      ^^
    ~\lstbox{\textSFxi{}}~
    = help: the index assignment already aliases the value implicitly, use [ instead
    ~\lstbox{\textSFxi{}}~
11  -     a:[0] = 0;
11  +     a[0] = 0;
    ~\lstbox{\textSFii{}}~~\lstbox{\textSFx{}}~~\lstbox{\textSFx{}}~~\lstbox{\textSFx{}}~~\lstbox{\textSFx{}}~~\lstbox{\textSFx{}}~~\lstbox{\textSFx{}}~~\lstbox{\textSFx{}}~
\end{lstlisting}

\subsubsection*{How to fix it}

Use the plain \tokennolst{[} index operator on the left of the assignment. The assignment already aliases the value it writes, which is the edit the diagnostic proposes.

\errorcode{E4014}{class alias operator :. is not usable to access fields}
\label{sec:codes:e4014}

Raised during \textbf{validation}, from \texttt{ValidateErrorMessage::\allowbreak{}CLASS\_\allowbreak{}ALIAS\_\allowbreak{}FIELD} (\textit{src/\allowbreak{}ymirc/\allowbreak{}semantic/\allowbreak{}validator/\allowbreak{}errors.yr}).

\subsubsection*{What it means}

The class alias operator \tokennolst{:.} was used to access a field. That operator produces a borrowed view for calling a method; a field access takes its mutability from the value itself, so the operator has nothing to do.

\subsubsection*{Example}

\textit{test\_\allowbreak{}resources/\allowbreak{}diagnostics/\allowbreak{}test98.yr}:

%% check: skip
\begin{lstlisting}[style=coloredverbatimError]
// the alias operator used to read a field of a map
fn main () {
    let mut m = copy [1 => 2];
    let _ = m:.len;
}
\end{lstlisting}

\begin{lstlisting}[style=bashVerb, breaklines=true, escapechar=~]
Error[E4014] : class alias operator :. is not usable to access fields
 --> test_resources/diagnostics/test98.yr:(4,14)
 4  ~\lstbox{\textSFxi{}}~     let _ = m:.len;
    ~\lstbox{\textSFv{}}~              ^^---
    ~\lstbox{\textSFxi{}}~
    = when validating test98::main -- (test_resources/diagnostics/test98.yr:(2,4))
    ~\lstbox{\textSFii{}}~~\lstbox{\textSFx{}}~~\lstbox{\textSFx{}}~~\lstbox{\textSFx{}}~~\lstbox{\textSFx{}}~~\lstbox{\textSFx{}}~~\lstbox{\textSFx{}}~~\lstbox{\textSFx{}}~
\end{lstlisting}

\subsubsection*{How to fix it}

Use a plain \tokennolst{.} to access the field.

\errorcode{E4015}{alias operator :. is only usable on class, record, entity, map or generator types, not \errhole{}}
\label{sec:codes:e4015}

Raised during \textbf{validation}, from \texttt{ValidateErrorMessage::\allowbreak{}CLASS\_\allowbreak{}ALIAS\_\allowbreak{}FIELD\_\allowbreak{}NO\_\allowbreak{}CLASS} (\textit{src/\allowbreak{}ymirc/\allowbreak{}semantic/\allowbreak{}validator/\allowbreak{}errors.yr}).

\subsubsection*{What it means}

The alias operator \tokennolst{:.} was applied to a type that has no method to alias -- it is only defined for classes, records, entities, maps and generators.

\subsubsection*{Example}

\textit{test\_\allowbreak{}resources/\allowbreak{}diagnostics/\allowbreak{}test23.yr}:

%% check: skip
\begin{lstlisting}[style=coloredverbatimError]
fn main () {
    let x = 1;
    x:.foo ();
}
\end{lstlisting}

\begin{lstlisting}[style=bashVerb, breaklines=true, escapechar=~]
Error[E4015] : alias operator :. is only usable on class, record, entity, map or generator types, not x
 --> test_resources/diagnostics/test23.yr:(3,5)
 3  ~\lstbox{\textSFxi{}}~     x:.foo ();
    ~\lstbox{\textSFv{}}~     ^
    ~\lstbox{\textSFii{}}~~\lstbox{\textSFx{}}~~\lstbox{\textSFx{}}~~\lstbox{\textSFx{}}~~\lstbox{\textSFx{}}~~\lstbox{\textSFx{}}~~\lstbox{\textSFx{}}~~\lstbox{\textSFx{}}~
\end{lstlisting}

\subsubsection*{How to fix it}

Drop the operator, or apply it to a value of one of those types.

\errorcode{E4016}{class alias operator :[ is only usable on class, record or entity types, not \errhole{}}
\label{sec:codes:e4016}

Raised during \textbf{validation}, from \texttt{ValidateErrorMessage::\allowbreak{}CLASS\_\allowbreak{}ALIAS\_\allowbreak{}INDEX\_\allowbreak{}NO\_\allowbreak{}CLASS} (\textit{src/\allowbreak{}ymirc/\allowbreak{}semantic/\allowbreak{}validator/\allowbreak{}errors.yr}).

\subsubsection*{What it means}

The class alias index operator \tokennolst{:[} was applied to a type it is not defined for; it only applies to classes, records and entities.

\subsubsection*{Example}

\textit{test\_\allowbreak{}resources/\allowbreak{}diagnostics/\allowbreak{}test24.yr}:

%% check: skip
\begin{lstlisting}[style=coloredverbatimError]
fn main () {
    let x = 1;
    x:[0];
}
\end{lstlisting}

\begin{lstlisting}[style=bashVerb, breaklines=true, escapechar=~]
Error[E4016] : class alias operator :[ is only usable on class, record or entity types, not x
 --> test_resources/diagnostics/test24.yr:(3,5)
 3  ~\lstbox{\textSFxi{}}~     x:[0];
    ~\lstbox{\textSFv{}}~     ^
    ~\lstbox{\textSFii{}}~~\lstbox{\textSFx{}}~~\lstbox{\textSFx{}}~~\lstbox{\textSFx{}}~~\lstbox{\textSFx{}}~~\lstbox{\textSFx{}}~~\lstbox{\textSFx{}}~~\lstbox{\textSFx{}}~
\end{lstlisting}

\subsubsection*{How to fix it}

Use the plain \tokennolst{[} index operator.

\errorcode{E4029}{movable type \errhole{} cannot be embedded within another type}
\label{sec:codes:e4029}

Raised during \textbf{validation}, from \texttt{ValidateErrorMessage::\allowbreak{}CONTAIN\_\allowbreak{}MOVABLE\_\allowbreak{}TYPE} (\textit{src/\allowbreak{}ymirc/\allowbreak{}semantic/\allowbreak{}validator/\allowbreak{}errors.yr}).

\subsubsection*{What it means}

A movable type was embedded in another type. A movable type has an identity that survives being moved, which embedding it by value would break.

\subsubsection*{Example}

\textit{test\_\allowbreak{}resources/\allowbreak{}union/\allowbreak{}test4.yr}:

%% check: skip
\begin{lstlisting}[style=coloredverbatimError]
entity Z {
    pub self () {}
}

@overlaid
record A {
    let z : Z;
}
\end{lstlisting}

\begin{lstlisting}[style=bashVerb, breaklines=true, escapechar=~]
Error[E4029] : movable type test4::Z cannot be embedded within another type
 --> test_resources/union/test4.yr:(7,5)
 7  ~\lstbox{\textSFxi{}}~     let z : Z;
    ~\lstbox{\textSFv{}}~     ^^^     ^
    ~\lstbox{\textSFxi{}}~ Note : the structure is defined as a record but must be an entity to store movable values
    ~\lstbox{\textSFxi{}}~  --> test_resources/union/test4.yr:(6,8)
    ~\lstbox{\textSFxi{}}~  6  ~\lstbox{\textSFxi{}}~ record A {
    ~\lstbox{\textSFxi{}}~     ~\lstbox{\textSFv{}}~        ^
    ~\lstbox{\textSFxi{}}~     ~\lstbox{\textSFii{}}~~\lstbox{\textSFx{}}~~\lstbox{\textSFx{}}~~\lstbox{\textSFx{}}~~\lstbox{\textSFx{}}~~\lstbox{\textSFx{}}~~\lstbox{\textSFx{}}~~\lstbox{\textSFx{}}~
    ~\lstbox{\textSFii{}}~~\lstbox{\textSFx{}}~~\lstbox{\textSFx{}}~~\lstbox{\textSFx{}}~~\lstbox{\textSFx{}}~~\lstbox{\textSFx{}}~~\lstbox{\textSFvii{}}~~\lstbox{\textSFx{}}~
\end{lstlisting}

\subsubsection*{How to fix it}

Hold it behind a class or a pointer, rather than by value.

\errorcode{E4045}{discarding the constant property is prohibited}
\label{sec:codes:e4045}

Raised during \textbf{validation}, from \texttt{ValidateErrorMessage::\allowbreak{}DISCARD\_\allowbreak{}CONST} (\textit{src/\allowbreak{}ymirc/\allowbreak{}semantic/\allowbreak{}validator/\allowbreak{}errors.yr}).

\subsubsection*{What it means}

An operation would drop the constant property of a value -- turning something borrowed as constant into something writable.

\subsubsection*{Example}

\textit{test\_\allowbreak{}resources/\allowbreak{}lit\_\allowbreak{}tuples/\allowbreak{}test13.yr}:

%% check: skip
\begin{lstlisting}[style=coloredverbatimError]
fn main () {
    let dmut a = (copy [1, 2],);
    let dmut _ = a.0;
}
\end{lstlisting}

\begin{lstlisting}[style=bashVerb, breaklines=true, escapechar=~]
Error[E4045] : discarding the constant property is prohibited
 --> test_resources/lit_tuples/test13.yr:(3,14)
 3  ~\lstbox{\textSFxi{}}~     let dmut _ = a.0;
    ~\lstbox{\textSFv{}}~              ^
    ~\lstbox{\textSFxi{}}~ Note : implicit alias of type mut [mut i32] is prohibited, it implicitly gives up on borrowings of mutable values
    ~\lstbox{\textSFxi{}}~  --> test_resources/lit_tuples/test13.yr:(3,19)
    ~\lstbox{\textSFxi{}}~  3  ~\lstbox{\textSFxi{}}~     let dmut _ = a.0;
    ~\lstbox{\textSFxi{}}~     ~\lstbox{\textSFv{}}~                   ^
    ~\lstbox{\textSFxi{}}~     ~\lstbox{\textSFii{}}~~\lstbox{\textSFx{}}~~\lstbox{\textSFx{}}~~\lstbox{\textSFx{}}~~\lstbox{\textSFx{}}~~\lstbox{\textSFx{}}~~\lstbox{\textSFx{}}~~\lstbox{\textSFx{}}~
    ~\lstbox{\textSFii{}}~~\lstbox{\textSFx{}}~~\lstbox{\textSFx{}}~~\lstbox{\textSFx{}}~~\lstbox{\textSFx{}}~~\lstbox{\textSFx{}}~~\lstbox{\textSFvii{}}~~\lstbox{\textSFx{}}~
\end{lstlisting}

\subsubsection*{How to fix it}

Take a copy if the value has to be modified, or keep the operation constant. The mutability of a borrowed value is part of its type and cannot be cast away.

\errorcode{E4082}{left operand of type \errhole{} is immutable}
\label{sec:codes:e4082}

Raised during \textbf{validation}, from \texttt{ValidateErrorMessage::\allowbreak{}IMMUTABLE\_\allowbreak{}LVALUE} (\textit{src/\allowbreak{}ymirc/\allowbreak{}semantic/\allowbreak{}validator/\allowbreak{}errors.yr}).

\subsubsection*{What it means}

An assignment writes through a value whose type is immutable.

\subsubsection*{Example}

\textit{test\_\allowbreak{}resources/\allowbreak{}for\_\allowbreak{}loops/\allowbreak{}tuples/\allowbreak{}test2.yr}:

%% check: skip
\begin{lstlisting}[style=coloredverbatimError]
fn main () {
    let a = (1, 2, 3);
    for x in ref a {
        x;
    }
}
\end{lstlisting}

\begin{lstlisting}[style=bashVerb, breaklines=true, escapechar=~]
Error[E4082] : left operand of type (i32, i32, i32) is immutable
 --> test_resources/for_loops/tuples/test2.yr:(3,14)
 3  ~\lstbox{\textSFxi{}}~     for x in ref a {
    ~\lstbox{\textSFv{}}~              ^^^ ^
    ~\lstbox{\textSFii{}}~~\lstbox{\textSFx{}}~~\lstbox{\textSFx{}}~~\lstbox{\textSFx{}}~~\lstbox{\textSFx{}}~~\lstbox{\textSFx{}}~~\lstbox{\textSFx{}}~~\lstbox{\textSFx{}}~
\end{lstlisting}

\subsubsection*{How to fix it}

Declare the variable, or the borrow it came from, as mutable. Mutability travels with the type, so a value reached through a constant borrow stays constant.

\errorcode{E4083}{type must be immutable}
\label{sec:codes:e4083}

Raised during \textbf{validation}, from \texttt{ValidateErrorMessage::\allowbreak{}IMMUTABLE\_\allowbreak{}PARENT} (\textit{src/\allowbreak{}ymirc/\allowbreak{}semantic/\allowbreak{}validator/\allowbreak{}errors.yr}).

\subsubsection*{What it means}

A type used in this position has to be immutable and is not -- typically the aliased type of a \tokennolst{def}, or a type used where sharing requires immutability.

\subsubsection*{Example}

\textit{test\_\allowbreak{}resources/\allowbreak{}aka/\allowbreak{}test2.yr}:

%% check: skip
\begin{lstlisting}[style=coloredverbatimError]
def A : dmut [i32];

fn main ()
{
    let dmut x : A = copy [1, 2, 3];
    x [0] = 1;
}
\end{lstlisting}

\begin{lstlisting}[style=bashVerb, breaklines=true, escapechar=~]
Error[E4083] : type must be immutable
 --> test_resources/aka/test2.yr:(1,9)
 1  ~\lstbox{\textSFxi{}}~ def A : dmut [i32];
    ~\lstbox{\textSFv{}}~         ^^^^
    ~\lstbox{\textSFii{}}~~\lstbox{\textSFx{}}~~\lstbox{\textSFx{}}~~\lstbox{\textSFx{}}~~\lstbox{\textSFx{}}~~\lstbox{\textSFx{}}~~\lstbox{\textSFx{}}~~\lstbox{\textSFx{}}~
\end{lstlisting}

\subsubsection*{How to fix it}

Remove the mutability decorator from the type.

\errorcode{E4088}{implicit move of type \errhole{} is prohibited}
\label{sec:codes:e4088}

Raised during \textbf{validation}, from \texttt{ValidateErrorMessage::\allowbreak{}IMPLICIT\_\allowbreak{}MOVE} (\textit{src/\allowbreak{}ymirc/\allowbreak{}semantic/\allowbreak{}validator/\allowbreak{}errors.yr}).

\subsubsection*{What it means}

A value would be moved implicitly. A move leaves its source unusable, so it is always written out.

\subsubsection*{Example}

\textit{test\_\allowbreak{}resources/\allowbreak{}diagnostics/\allowbreak{}test106.yr}:

%% check: skip
\begin{lstlisting}[style=coloredverbatimError]
// an entity bound to a second variable without a move
entity Counter {
    pub self () {}
}

fn main () {
    let c = Counter ();
    let _d_ = c;
}
\end{lstlisting}

\begin{lstlisting}[style=bashVerb, breaklines=true, escapechar=~]
Error[E4088] : implicit move of type test106::Counter is prohibited
 --> test_resources/diagnostics/test106.yr:(8,9)
 8  ~\lstbox{\textSFxi{}}~     let _d_ = c;
    ~\lstbox{\textSFv{}}~         ^^^   -
    ~\lstbox{\textSFxi{}}~
    = when validating test106::main -- (test_resources/diagnostics/test106.yr:(6,4))
    ~\lstbox{\textSFii{}}~~\lstbox{\textSFx{}}~~\lstbox{\textSFx{}}~~\lstbox{\textSFx{}}~~\lstbox{\textSFx{}}~~\lstbox{\textSFx{}}~~\lstbox{\textSFx{}}~~\lstbox{\textSFx{}}~
\end{lstlisting}

\subsubsection*{How to fix it}

Write \tokennolst{move} before the value, or pass a copy instead.

\errorcode{E4089}{ref of type \errhole{} must be mutable, but is const}
\label{sec:codes:e4089}

Raised during \textbf{validation}, from \texttt{ValidateErrorMessage::\allowbreak{}IMPLICIT\_\allowbreak{}MUT\_\allowbreak{}REFERENCE} (\textit{src/\allowbreak{}ymirc/\allowbreak{}semantic/\allowbreak{}validator/\allowbreak{}errors.yr}).

\subsubsection*{What it means}

A \tokennolst{ref} parameter is declared mutable and a constant value was passed.

\subsubsection*{Example}

\textit{test\_\allowbreak{}resources/\allowbreak{}for\_\allowbreak{}loops/\allowbreak{}tuples/\allowbreak{}test5.yr}:

%% check: skip
\begin{lstlisting}[style=coloredverbatimError]
fn main () {
    let dmut a = (copy [1, 2],);
    for ref mut x in a {
        x;
    }
}
\end{lstlisting}

\begin{lstlisting}[style=bashVerb, breaklines=true, escapechar=~]
Error[E4089] : ref of type mut [mut i32] must be mutable, but is const
 --> test_resources/for_loops/tuples/test5.yr:(3,17)
 3  ~\lstbox{\textSFxi{}}~     for ref mut x in a {
    ~\lstbox{\textSFv{}}~                 ^
    ~\lstbox{\textSFii{}}~~\lstbox{\textSFx{}}~~\lstbox{\textSFx{}}~~\lstbox{\textSFx{}}~~\lstbox{\textSFx{}}~~\lstbox{\textSFx{}}~~\lstbox{\textSFx{}}~~\lstbox{\textSFx{}}~
\end{lstlisting}

\subsubsection*{How to fix it}

Pass a mutable value, or declare the parameter constant.

\errorcode{E4092}{cannot pass value of type \errhole{} as ref}
\label{sec:codes:e4092}

Raised during \textbf{validation}, from \texttt{ValidateErrorMessage::\allowbreak{}IMPLICIT\_\allowbreak{}REFERENCE} (\textit{src/\allowbreak{}ymirc/\allowbreak{}semantic/\allowbreak{}validator/\allowbreak{}errors.yr}).

\subsubsection*{What it means}

A value was passed as \tokennolst{ref} that has no storage to borrow -- a temporary, or the result of an expression.

\subsubsection*{Example}

\textit{test\_\allowbreak{}resources/\allowbreak{}control\_\allowbreak{}flow/\allowbreak{}test21.yr}:

%% check: skip
\begin{lstlisting}[style=coloredverbatimError]
fn main () {
    let (ref x,) = (1,);
}
\end{lstlisting}

\begin{lstlisting}[style=bashVerb, breaklines=true, escapechar=~]
Error[E4092] : cannot pass value of type i32 as ref
 --> test_resources/control_flow/test21.yr:(2,14)
 2  ~\lstbox{\textSFxi{}}~     let (ref x,) = (1,);
    ~\lstbox{\textSFv{}}~              ^      ^
    ~\lstbox{\textSFxi{}}~ Note : for pattern field access 0 on type (mut i32,)
    ~\lstbox{\textSFxi{}}~  --> test_resources/control_flow/test21.yr:(2,9)
    ~\lstbox{\textSFxi{}}~  2  ~\lstbox{\textSFxi{}}~     let (ref x,) = (1,);
    ~\lstbox{\textSFxi{}}~     ~\lstbox{\textSFv{}}~         ^
    ~\lstbox{\textSFxi{}}~     ~\lstbox{\textSFii{}}~~\lstbox{\textSFx{}}~~\lstbox{\textSFx{}}~~\lstbox{\textSFx{}}~~\lstbox{\textSFx{}}~~\lstbox{\textSFx{}}~~\lstbox{\textSFx{}}~~\lstbox{\textSFx{}}~
    ~\lstbox{\textSFii{}}~~\lstbox{\textSFx{}}~~\lstbox{\textSFx{}}~~\lstbox{\textSFx{}}~~\lstbox{\textSFx{}}~~\lstbox{\textSFx{}}~~\lstbox{\textSFvii{}}~~\lstbox{\textSFx{}}~
\end{lstlisting}

\subsubsection*{How to fix it}

Bind the value to a variable first and pass that.

\errorcode{E4099}{index assignment of a value of type \errhole{} requires a mutable value}
\label{sec:codes:e4099}

Raised during \textbf{validation}, from \texttt{ValidateErrorMessage::\allowbreak{}INDEX\_\allowbreak{}ASSIGN\_\allowbreak{}NOT\_\allowbreak{}MUTABLE} (\textit{src/\allowbreak{}ymirc/\allowbreak{}semantic/\allowbreak{}validator/\allowbreak{}errors.yr}).

\subsubsection*{What it means}

An index assignment writes into a value that is not mutable.

\subsubsection*{Example}

\textit{test\_\allowbreak{}resources/\allowbreak{}lit\_\allowbreak{}class/\allowbreak{}operators/\allowbreak{}test17.yr}:

%% check: skip
\begin{lstlisting}[style=coloredverbatimError]
class A {
    pub self () {}

    pub fn opIndexAssign (mut self, _ : i32, _ : i32) {}

}


fn main () {
    let dmut a = copy A ();
    a [0] = 0;
    a:[0] = 0;

    let b = copy A ();
    b [0] = 0;
    b:[0] = 0;
}
\end{lstlisting}

\begin{lstlisting}[style=bashVerb, breaklines=true, escapechar=~]
Error[E4099] : index assignment of a value of type &(test17::A) requires a mutable value
 --> test_resources/lit_class/operators/test17.yr:(15,5)
15  ~\lstbox{\textSFxi{}}~     b [0] = 0;
    ~\lstbox{\textSFv{}}~     ^ ^  note: the index assignment calls the mutable method opIndexAssign, and thus implicitly aliases the value
    ~\lstbox{\textSFii{}}~~\lstbox{\textSFx{}}~~\lstbox{\textSFx{}}~~\lstbox{\textSFx{}}~~\lstbox{\textSFx{}}~~\lstbox{\textSFx{}}~~\lstbox{\textSFx{}}~~\lstbox{\textSFx{}}~
\end{lstlisting}

\subsubsection*{How to fix it}

Declare the indexed value mutable, cf. \hyperref[sec:codes:e4082]{E4082}.

\errorcode{E4125}{cannot alias an expand value, maybe alias and expand keywords are inverted?}
\label{sec:codes:e4125}

Raised during \textbf{validation}, from \texttt{ValidateErrorMessage::\allowbreak{}MISMATCH\_\allowbreak{}ALIAS\_\allowbreak{}EXPAND} (\textit{src/\allowbreak{}ymirc/\allowbreak{}semantic/\allowbreak{}validator/\allowbreak{}errors.yr}).

\subsubsection*{What it means}

\tokennolst{alias} was applied to an \tokennolst{expand} value. The two keywords are in the order that aliases the expansion rather than expanding the alias, which is almost never what is meant.

\subsubsection*{Example}

\textit{test\_\allowbreak{}resources/\allowbreak{}diagnostics/\allowbreak{}test138.yr}:

%% check: skip
\begin{lstlisting}[style=coloredverbatimError]
// alias and expand written the wrong way round
fn take (dmut _ : [i32], _ : [i32]) {}

fn main () {
    let mut z : (mut [mut i32], [i32]) = alias (copy [1, 2], copy [3, 4]);
    take (alias expand z);
}
\end{lstlisting}

\begin{lstlisting}[style=bashVerb, breaklines=true, escapechar=~]
Error[E4125] : cannot alias an expand value, maybe alias and expand keywords are inverted?
 --> test_resources/diagnostics/test138.yr:(6,11)
 6  ~\lstbox{\textSFxi{}}~     take (alias expand z);
    ~\lstbox{\textSFv{}}~           ^^^^^ ------
    ~\lstbox{\textSFxi{}}~
    = help: invert the two keywords
    ~\lstbox{\textSFxi{}}~
 6  -     take (alias expand z);
 6  +     take (expand alias z);
    ~\lstbox{\textSFxi{}}~
    = when validating test138::main -- (test_resources/diagnostics/test138.yr:(4,4))
    ~\lstbox{\textSFii{}}~~\lstbox{\textSFx{}}~~\lstbox{\textSFx{}}~~\lstbox{\textSFx{}}~~\lstbox{\textSFx{}}~~\lstbox{\textSFx{}}~~\lstbox{\textSFx{}}~~\lstbox{\textSFx{}}~
\end{lstlisting}

\subsubsection*{How to fix it}

Swap them: \tokennolst{expand alias v}.

\errorcode{E4126}{cannot alias a lazy value, maybe alias and lazy keywords are inverted?}
\label{sec:codes:e4126}

Raised during \textbf{validation}, from \texttt{ValidateErrorMessage::\allowbreak{}MISMATCH\_\allowbreak{}ALIAS\_\allowbreak{}LAZY} (\textit{src/\allowbreak{}ymirc/\allowbreak{}semantic/\allowbreak{}validator/\allowbreak{}errors.yr}).

\subsubsection*{What it means}

\tokennolst{alias} was applied to a \tokennolst{lazy} value; the keywords are inverted.

\subsubsection*{Example}

\textit{test\_\allowbreak{}resources/\allowbreak{}lazyness/\allowbreak{}test6.yr}:

%% check: skip
\begin{lstlisting}[style=coloredverbatimError]
fn main () {
    let dmut i = copy [1, 2, 3];
    let lazy dmut a = copy (lazy i);
    let lazy dmut b = alias (lazy i);
}
\end{lstlisting}

\begin{lstlisting}[style=bashVerb, breaklines=true, escapechar=~]
Error[E4126] : cannot alias a lazy value, maybe alias and lazy keywords are inverted?
 --> test_resources/lazyness/test6.yr:(4,23)
 4  ~\lstbox{\textSFxi{}}~     let lazy dmut b = alias (lazy i);
    ~\lstbox{\textSFv{}}~                       ^^^^^  ----
    ~\lstbox{\textSFxi{}}~
    = help: invert the two keywords
    ~\lstbox{\textSFxi{}}~
 4  -     let lazy dmut b = alias (lazy i);
 4  +     let lazy dmut b = lazy (alias i);
    ~\lstbox{\textSFii{}}~~\lstbox{\textSFx{}}~~\lstbox{\textSFx{}}~~\lstbox{\textSFx{}}~~\lstbox{\textSFx{}}~~\lstbox{\textSFx{}}~~\lstbox{\textSFx{}}~~\lstbox{\textSFx{}}~
\end{lstlisting}

\subsubsection*{How to fix it}

Swap them: \tokennolst{lazy alias v}.

\errorcode{E4127}{cannot copy an expand value, maybe copy and expand keywords are inverted?}
\label{sec:codes:e4127}

Raised during \textbf{validation}, from \texttt{ValidateErrorMessage::\allowbreak{}MISMATCH\_\allowbreak{}COPY\_\allowbreak{}EXPAND} (\textit{src/\allowbreak{}ymirc/\allowbreak{}semantic/\allowbreak{}validator/\allowbreak{}errors.yr}).

\subsubsection*{What it means}

\tokennolst{copy} was applied to an \tokennolst{expand} value; the keywords are inverted.

\subsubsection*{Example}

\textit{test\_\allowbreak{}resources/\allowbreak{}diagnostics/\allowbreak{}test139.yr}:

%% check: skip
\begin{lstlisting}[style=coloredverbatimError]
// copy and expand written the wrong way round
fn take (_ : i32, _ : i32) {}

fn main () {
    let x = [1, 2];
    take (copy expand x);
}
\end{lstlisting}

\begin{lstlisting}[style=bashVerb, breaklines=true, escapechar=~]
Error[E4127] : cannot copy an expand value, maybe copy and expand keywords are inverted?
 --> test_resources/diagnostics/test139.yr:(6,11)
 6  ~\lstbox{\textSFxi{}}~     take (copy expand x);
    ~\lstbox{\textSFv{}}~           ^^^^ ------
    ~\lstbox{\textSFxi{}}~
    = help: invert the two keywords
    ~\lstbox{\textSFxi{}}~
 6  -     take (copy expand x);
 6  +     take (expand copy x);
    ~\lstbox{\textSFxi{}}~
    = when validating test139::main -- (test_resources/diagnostics/test139.yr:(4,4))
    ~\lstbox{\textSFii{}}~~\lstbox{\textSFx{}}~~\lstbox{\textSFx{}}~~\lstbox{\textSFx{}}~~\lstbox{\textSFx{}}~~\lstbox{\textSFx{}}~~\lstbox{\textSFx{}}~~\lstbox{\textSFx{}}~
\end{lstlisting}

\subsubsection*{How to fix it}

Swap them: \tokennolst{expand copy v}.

\errorcode{E4128}{cannot copy a lazy value, maybe copy and lazy keywords are inverted?}
\label{sec:codes:e4128}

Raised during \textbf{validation}, from \texttt{ValidateErrorMessage::\allowbreak{}MISMATCH\_\allowbreak{}COPY\_\allowbreak{}LAZY} (\textit{src/\allowbreak{}ymirc/\allowbreak{}semantic/\allowbreak{}validator/\allowbreak{}errors.yr}).

\subsubsection*{What it means}

\tokennolst{copy} was applied to a \tokennolst{lazy} value; the keywords are inverted.

\subsubsection*{Example}

\textit{test\_\allowbreak{}resources/\allowbreak{}lazyness/\allowbreak{}test6.yr}:

%% check: skip
\begin{lstlisting}[style=coloredverbatimError]
fn main () {
    let dmut i = copy [1, 2, 3];
    let lazy dmut a = copy (lazy i);
    let lazy dmut b = alias (lazy i);
}
\end{lstlisting}

\begin{lstlisting}[style=bashVerb, breaklines=true, escapechar=~]
Error[E4128] : cannot copy a lazy value, maybe copy and lazy keywords are inverted?
 --> test_resources/lazyness/test6.yr:(3,23)
 3  ~\lstbox{\textSFxi{}}~     let lazy dmut a = copy (lazy i);
    ~\lstbox{\textSFv{}}~                       ^^^^  ----
    ~\lstbox{\textSFxi{}}~
    = help: invert the two keywords
    ~\lstbox{\textSFxi{}}~
 3  -     let lazy dmut a = copy (lazy i);
 3  +     let lazy dmut a = lazy (copy i);
    ~\lstbox{\textSFii{}}~~\lstbox{\textSFx{}}~~\lstbox{\textSFx{}}~~\lstbox{\textSFx{}}~~\lstbox{\textSFx{}}~~\lstbox{\textSFx{}}~~\lstbox{\textSFx{}}~~\lstbox{\textSFx{}}~
\end{lstlisting}

\subsubsection*{How to fix it}

Swap them: \tokennolst{lazy copy v}.

\errorcode{E4129}{cannot copy an expand value, maybe dcopy and expand keywords are inverted?}
\label{sec:codes:e4129}

Raised during \textbf{validation}, from \texttt{ValidateErrorMessage::\allowbreak{}MISMATCH\_\allowbreak{}DCOPY\_\allowbreak{}EXPAND} (\textit{src/\allowbreak{}ymirc/\allowbreak{}semantic/\allowbreak{}validator/\allowbreak{}errors.yr}).

\subsubsection*{What it means}

\tokennolst{dcopy} was applied to an \tokennolst{expand} value; the keywords are inverted.

\subsubsection*{Example}

\textit{test\_\allowbreak{}resources/\allowbreak{}diagnostics/\allowbreak{}test140.yr}:

%% check: skip
\begin{lstlisting}[style=coloredverbatimError]
// dcopy and expand written the wrong way round
fn take (_ : i32, _ : i32) {}

fn main () {
    let x = [1, 2];
    take (dcopy expand x);
}
\end{lstlisting}

\begin{lstlisting}[style=bashVerb, breaklines=true, escapechar=~]
Error[E4129] : cannot copy an expand value, maybe dcopy and expand keywords are inverted?
 --> test_resources/diagnostics/test140.yr:(6,11)
 6  ~\lstbox{\textSFxi{}}~     take (dcopy expand x);
    ~\lstbox{\textSFv{}}~           ^^^^^ ------
    ~\lstbox{\textSFxi{}}~
    = help: invert the two keywords
    ~\lstbox{\textSFxi{}}~
 6  -     take (dcopy expand x);
 6  +     take (expand dcopy x);
    ~\lstbox{\textSFxi{}}~
    = when validating test140::main -- (test_resources/diagnostics/test140.yr:(4,4))
    ~\lstbox{\textSFii{}}~~\lstbox{\textSFx{}}~~\lstbox{\textSFx{}}~~\lstbox{\textSFx{}}~~\lstbox{\textSFx{}}~~\lstbox{\textSFx{}}~~\lstbox{\textSFx{}}~~\lstbox{\textSFx{}}~
\end{lstlisting}

\subsubsection*{How to fix it}

Swap them: \tokennolst{expand dcopy v}.

\errorcode{E4130}{cannot copy a lazy value, maybe dcopy and lazy keywords are inverted?}
\label{sec:codes:e4130}

Raised during \textbf{validation}, from \texttt{ValidateErrorMessage::\allowbreak{}MISMATCH\_\allowbreak{}DCOPY\_\allowbreak{}LAZY} (\textit{src/\allowbreak{}ymirc/\allowbreak{}semantic/\allowbreak{}validator/\allowbreak{}errors.yr}).

\subsubsection*{What it means}

\tokennolst{dcopy} was applied to a \tokennolst{lazy} value; the keywords are inverted.

\subsubsection*{Example}

\textit{test\_\allowbreak{}resources/\allowbreak{}lazyness/\allowbreak{}test19.yr}:

%% check: skip
\begin{lstlisting}[style=coloredverbatimError]
// a deep copy of a lazy value reads the wrong way round, the same way a
// shallow one does
fn main () {
    let dmut i = copy [1, 2, 3];
    let lazy dmut a = dcopy (lazy i);
    a;
}
\end{lstlisting}

\begin{lstlisting}[style=bashVerb, breaklines=true, escapechar=~]
Error[E4130] : cannot copy a lazy value, maybe dcopy and lazy keywords are inverted?
 --> test_resources/lazyness/test19.yr:(5,23)
 5  ~\lstbox{\textSFxi{}}~     let lazy dmut a = dcopy (lazy i);
    ~\lstbox{\textSFv{}}~                       ^^^^^  ----
    ~\lstbox{\textSFxi{}}~
    = help: invert the two keywords
    ~\lstbox{\textSFxi{}}~
 5  -     let lazy dmut a = dcopy (lazy i);
 5  +     let lazy dmut a = lazy (dcopy i);
    ~\lstbox{\textSFxi{}}~
    = when validating test19::main -- (test_resources/lazyness/test19.yr:(3,4))
    ~\lstbox{\textSFii{}}~~\lstbox{\textSFx{}}~~\lstbox{\textSFx{}}~~\lstbox{\textSFx{}}~~\lstbox{\textSFx{}}~~\lstbox{\textSFx{}}~~\lstbox{\textSFx{}}~~\lstbox{\textSFx{}}~
\end{lstlisting}

\subsubsection*{How to fix it}

Swap them: \tokennolst{lazy dcopy v}.

\errorcode{E4137}{an iterator cannot be mutable, if it is not a reference or does not borrow mutable data}
\label{sec:codes:e4137}

Raised during \textbf{validation}, from \texttt{ValidateErrorMessage::\allowbreak{}MUTABLE\_\allowbreak{}CONST\_\allowbreak{}ITERATOR} (\textit{src/\allowbreak{}ymirc/\allowbreak{}semantic/\allowbreak{}validator/\allowbreak{}errors.yr}).

\subsubsection*{What it means}

A loop iterator is declared mutable while it neither is a reference nor borrows mutable data. Mutating it would change a copy, which the loop would then discard.

\subsubsection*{Example}

\textit{test\_\allowbreak{}resources/\allowbreak{}for\_\allowbreak{}loops/\allowbreak{}arrays/\allowbreak{}test2.yr}:

%% check: skip
\begin{lstlisting}[style=coloredverbatimError]
in test2;

fn main () {
    let dmut a = [1, 2, 3, 4];
    for mut elem : i32 in a {
        assert (a[i] == elem);
    }
}
\end{lstlisting}

\begin{lstlisting}[style=bashVerb, breaklines=true, escapechar=~]
Error[E4137] : an iterator cannot be mutable, if it is not a reference or does not borrow mutable data
 --> test_resources/for_loops/arrays/test2.yr:(5,9)
 5  ~\lstbox{\textSFxi{}}~     for mut elem : i32 in a {
    ~\lstbox{\textSFv{}}~         ^^^
    ~\lstbox{\textSFii{}}~~\lstbox{\textSFx{}}~~\lstbox{\textSFx{}}~~\lstbox{\textSFx{}}~~\lstbox{\textSFx{}}~~\lstbox{\textSFx{}}~~\lstbox{\textSFx{}}~~\lstbox{\textSFx{}}~
\end{lstlisting}

\subsubsection*{How to fix it}

Drop the \tokennolst{mut}, or iterate by reference if the elements are to be modified.

\errorcode{E4138}{a parameter cannot be mutable, if it is not a reference or does not borrow mutable data}
\label{sec:codes:e4138}

Raised during \textbf{validation}, from \texttt{ValidateErrorMessage::\allowbreak{}MUTABLE\_\allowbreak{}CONST\_\allowbreak{}PARAM} (\textit{src/\allowbreak{}ymirc/\allowbreak{}semantic/\allowbreak{}validator/\allowbreak{}errors.yr}).

\subsubsection*{What it means}

A parameter is declared mutable while it neither is a reference nor borrows mutable data, so the mutation would be invisible to the caller.

\subsubsection*{Example}

\textit{test\_\allowbreak{}resources/\allowbreak{}lazyness/\allowbreak{}test11.yr}:

%% check: skip
\begin{lstlisting}[style=coloredverbatimError]
fn foo (lazy mut a : i32) {}
\end{lstlisting}

\begin{lstlisting}[style=bashVerb, breaklines=true, escapechar=~]
Error[E4138] : a parameter cannot be mutable, if it is not a reference or does not borrow mutable data
 --> test_resources/lazyness/test11.yr:(1,14)
 1  ~\lstbox{\textSFxi{}}~ fn foo (lazy mut a : i32) {}
    ~\lstbox{\textSFv{}}~              ^^^
    ~\lstbox{\textSFii{}}~~\lstbox{\textSFx{}}~~\lstbox{\textSFx{}}~~\lstbox{\textSFx{}}~~\lstbox{\textSFx{}}~~\lstbox{\textSFx{}}~~\lstbox{\textSFx{}}~~\lstbox{\textSFx{}}~
\end{lstlisting}

\subsubsection*{How to fix it}

Drop the \tokennolst{mut}, or declare the parameter \tokennolst{ref} if the caller's value is to be modified.

\errorcode{E4140}{a lazy variable cannot be mutable, if it does not borrow mutable data}
\label{sec:codes:e4140}

Raised during \textbf{validation}, from \texttt{ValidateErrorMessage::\allowbreak{}MUTABLE\_\allowbreak{}LAZY\_\allowbreak{}VAR} (\textit{src/\allowbreak{}ymirc/\allowbreak{}semantic/\allowbreak{}validator/\allowbreak{}errors.yr}).

\subsubsection*{What it means}

A lazy binding is declared mutable while it does not borrow mutable data. It names a computation rather than a location: each read produces the value again in a fresh slot, so there is nothing a write could persist to.

A lazy \textit{global} is not in that case -- its value is stored inside the global itself and only computed on first read -- so \tokennolst{lazy mut A : i32 = 12;} is accepted and \tokennolst{A = 42;} writes to that storage.

\subsubsection*{Example}

\textit{test\_\allowbreak{}resources/\allowbreak{}lazyness/\allowbreak{}test20.yr}:

%% check: skip
\begin{lstlisting}[style=coloredverbatimError]
fn main () {
    let lazy mut a = lazy 1;
    a;
}
\end{lstlisting}

\begin{lstlisting}[style=bashVerb, breaklines=true, escapechar=~]
Error[E4140] : a lazy variable cannot be mutable, if it does not borrow mutable data
 --> test_resources/lazyness/test20.yr:(3,14)
 3  ~\lstbox{\textSFxi{}}~     let lazy mut a = lazy 1;
    ~\lstbox{\textSFv{}}~              ^^^
    ~\lstbox{\textSFii{}}~~\lstbox{\textSFx{}}~~\lstbox{\textSFx{}}~~\lstbox{\textSFx{}}~~\lstbox{\textSFx{}}~~\lstbox{\textSFx{}}~~\lstbox{\textSFx{}}~~\lstbox{\textSFx{}}~
\end{lstlisting}

\subsubsection*{How to fix it}

Drop the \tokennolst{mut}, or bind the value to an ordinary variable.

\errorcode{E4146}{\errhole{} is not an aliasable type}
\label{sec:codes:e4146}

Raised during \textbf{validation}, from \texttt{ValidateErrorMessage::\allowbreak{}NOT\_\allowbreak{}ALIASABLE} (\textit{src/\allowbreak{}ymirc/\allowbreak{}semantic/\allowbreak{}validator/\allowbreak{}errors.yr}).

\subsubsection*{What it means}

\tokennolst{alias} was applied to a type that does not borrow anything, so there is nothing to alias.

\subsubsection*{Example}

\textit{test\_\allowbreak{}resources/\allowbreak{}templates/\allowbreak{}implicit/\allowbreak{}test16.yr}:

%% check: skip
\begin{lstlisting}[style=coloredverbatimError]
fn foo {alias T} (a : T) {
    a;
}

fn main () {
    foo ([1, 2, 3]);
    foo ((1, 2, 3));
}
\end{lstlisting}

\begin{lstlisting}[style=bashVerb, breaklines=true, escapechar=~]
Error[E4146] : [i32 ; 3us] is not an aliasable type
 --> test_resources/templates/implicit/test16.yr:(1,15)
 1  ~\lstbox{\textSFxi{}}~ fn foo {alias T} (a : T) {
    ~\lstbox{\textSFv{}}~               ^
   ...
 --> test_resources/templates/implicit/test16.yr:(6,10)
 6  ~\lstbox{\textSFxi{}}~     foo ([1, 2, 3]);
    ~\lstbox{\textSFv{}}~          ^
    ~\lstbox{\textSFii{}}~~\lstbox{\textSFx{}}~~\lstbox{\textSFx{}}~~\lstbox{\textSFx{}}~~\lstbox{\textSFx{}}~~\lstbox{\textSFx{}}~~\lstbox{\textSFx{}}~~\lstbox{\textSFx{}}~
\end{lstlisting}

\subsubsection*{How to fix it}

Remove the \tokennolst{alias}.

\errorcode{E4151}{not a lvalue}
\label{sec:codes:e4151}

Raised during \textbf{validation}, from \texttt{ValidateErrorMessage::\allowbreak{}NOT\_\allowbreak{}A\_\allowbreak{}LVALUE} (\textit{src/\allowbreak{}ymirc/\allowbreak{}semantic/\allowbreak{}validator/\allowbreak{}errors.yr}).

\subsubsection*{What it means}

An expression was used where storage is required -- the left of an assignment, the target of a \tokennolst{ref}, the operand of an address-of -- and it has no storage of its own.

\subsubsection*{Example}

\textit{test\_\allowbreak{}resources/\allowbreak{}lit\_\allowbreak{}tuples/\allowbreak{}test20.yr}:

%% check: skip
\begin{lstlisting}[style=coloredverbatimError]
fn foo ()-> (i32, i32);

fn main () {
    let mut a = 0;

    (a, 1) = foo ();
}
\end{lstlisting}

\begin{lstlisting}[style=bashVerb, breaklines=true, escapechar=~]
Error[E4151] : not a lvalue
 --> test_resources/lit_tuples/test20.yr:(6,9)
 6  ~\lstbox{\textSFxi{}}~     (a, 1) = foo ();
    ~\lstbox{\textSFv{}}~         ^  ^
    ~\lstbox{\textSFii{}}~~\lstbox{\textSFx{}}~~\lstbox{\textSFx{}}~~\lstbox{\textSFx{}}~~\lstbox{\textSFx{}}~~\lstbox{\textSFx{}}~~\lstbox{\textSFx{}}~~\lstbox{\textSFx{}}~
\end{lstlisting}

\subsubsection*{How to fix it}

Bind the value to a variable and use that. The more specific diagnostics E4152 to E4157 say why a particular value has no storage.

\errorcode{E4153}{\errhole{} is a value iterator, and cannot be used as a lvalue}
\label{sec:codes:e4153}

Raised during \textbf{validation}, from \texttt{ValidateErrorMessage::\allowbreak{}NOT\_\allowbreak{}A\_\allowbreak{}LVALUE\_\allowbreak{}ITERATOR} (\textit{src/\allowbreak{}ymirc/\allowbreak{}semantic/\allowbreak{}validator/\allowbreak{}errors.yr}).

\subsubsection*{What it means}

A \tokennolst{for} loop iterator that yields values, not references, was used as an lvalue. Each iteration produces a fresh value, which is discarded at the end of the round.

\subsubsection*{Example}

\textit{test\_\allowbreak{}resources/\allowbreak{}for\_\allowbreak{}loops/\allowbreak{}tuples/\allowbreak{}test9.yr}:

%% check: skip
\begin{lstlisting}[style=coloredverbatimError]
fn main () {
    let dmut a = (copy [1, 2],);
    for dmut x in alias a {
        x = copy [8, 3];
    }
}
\end{lstlisting}

\begin{lstlisting}[style=bashVerb, breaklines=true, escapechar=~]
Error[E4153] : x is a value iterator, and cannot be used as a lvalue
 --> test_resources/for_loops/tuples/test9.yr:(4,9)
 4  ~\lstbox{\textSFxi{}}~         x = copy [8, 3];
    ~\lstbox{\textSFv{}}~         ^ ^
    ~\lstbox{\textSFii{}}~~\lstbox{\textSFx{}}~~\lstbox{\textSFx{}}~~\lstbox{\textSFx{}}~~\lstbox{\textSFx{}}~~\lstbox{\textSFx{}}~~\lstbox{\textSFx{}}~~\lstbox{\textSFx{}}~
\end{lstlisting}

\subsubsection*{How to fix it}

Iterate by reference so the loop yields the elements themselves, or collect the new values into a result.

\errorcode{E4154}{\errhole{} is a lazy variable, and cannot be used as a lvalue}
\label{sec:codes:e4154}

Raised during \textbf{validation}, from \texttt{ValidateErrorMessage::\allowbreak{}NOT\_\allowbreak{}A\_\allowbreak{}LVALUE\_\allowbreak{}LAZY} (\textit{src/\allowbreak{}ymirc/\allowbreak{}semantic/\allowbreak{}validator/\allowbreak{}errors.yr}).

\subsubsection*{What it means}

A lazy parameter or binding was used as an lvalue. It names a computation, not a location: each read produces the value again in a fresh slot, so there is no storage a write could persist to. A lazy global is not concerned -- its value is stored inside the global itself -- and a mutable one can be assigned.

\subsubsection*{Example}

\textit{test\_\allowbreak{}resources/\allowbreak{}lazyness/\allowbreak{}test7.yr}:

%% check: skip
\begin{lstlisting}[style=coloredverbatimError]
fn main () {
    let lazy dmut a = lazy copy [1, 2, 3];

    a = copy [1, 2, 3];
    a = lazy copy [2, 3, 4];
}
\end{lstlisting}

\begin{lstlisting}[style=bashVerb, breaklines=true, escapechar=~]
Error[E4154] : a is a lazy variable, and cannot be used as a lvalue
 --> test_resources/lazyness/test7.yr:(4,5)
 4  ~\lstbox{\textSFxi{}}~     a = copy [1, 2, 3];
    ~\lstbox{\textSFv{}}~     ^ ^
    ~\lstbox{\textSFii{}}~~\lstbox{\textSFx{}}~~\lstbox{\textSFx{}}~~\lstbox{\textSFx{}}~~\lstbox{\textSFx{}}~~\lstbox{\textSFx{}}~~\lstbox{\textSFx{}}~~\lstbox{\textSFx{}}~
\end{lstlisting}

\subsubsection*{How to fix it}

Bind the lazy value to an ordinary variable first.

\errorcode{E4155}{\errhole{} is a map access, and cannot be used as a lvalue}
\label{sec:codes:e4155}

Raised during \textbf{validation}, from \texttt{ValidateErrorMessage::\allowbreak{}NOT\_\allowbreak{}A\_\allowbreak{}LVALUE\_\allowbreak{}MAP\_\allowbreak{}ACCESS} (\textit{src/\allowbreak{}ymirc/\allowbreak{}semantic/\allowbreak{}validator/\allowbreak{}errors.yr}).

\subsubsection*{What it means}

A map access was used as an lvalue. A map lookup produces the value stored, not the slot it lives in, so writing to it would not reach the map.

\subsubsection*{Example}

\textit{test\_\allowbreak{}resources/\allowbreak{}lit\_\allowbreak{}map/\allowbreak{}test10.yr}:

%% check: skip
\begin{lstlisting}[style=coloredverbatimError]
fn main () {
    let dmut a : [[c8] => [i32]] = copy [];

    a ["foo"] = copy [0];
    if let Ok (ref dmut z) = ref a ["foo"] {
        z ~= [8];
    }
}
\end{lstlisting}

\begin{lstlisting}[style=bashVerb, breaklines=true, escapechar=~]
Error[E4155] : a ["foo"s8] is a map access, and cannot be used as a lvalue
 --> test_resources/lit_map/test10.yr:(5,30)
 5  ~\lstbox{\textSFxi{}}~     if let Ok (ref dmut z) = ref a ["foo"] {
    ~\lstbox{\textSFv{}}~                              ^^^   ^
    ~\lstbox{\textSFii{}}~~\lstbox{\textSFx{}}~~\lstbox{\textSFx{}}~~\lstbox{\textSFx{}}~~\lstbox{\textSFx{}}~~\lstbox{\textSFx{}}~~\lstbox{\textSFx{}}~~\lstbox{\textSFx{}}~
\end{lstlisting}

\subsubsection*{How to fix it}

Use the map's insertion operation to write, rather than assigning to the lookup.

\errorcode{E4156}{\errhole{} is a value parameter, and cannot be used as a lvalue}
\label{sec:codes:e4156}

Raised during \textbf{validation}, from \texttt{ValidateErrorMessage::\allowbreak{}NOT\_\allowbreak{}A\_\allowbreak{}LVALUE\_\allowbreak{}PARAM} (\textit{src/\allowbreak{}ymirc/\allowbreak{}semantic/\allowbreak{}validator/\allowbreak{}errors.yr}).

\subsubsection*{What it means}

A parameter passed by value was used as an lvalue in a way that expects the caller's storage.

\subsubsection*{Example}

\textit{test\_\allowbreak{}resources/\allowbreak{}diagnostics/\allowbreak{}test59.yr}:

%% check: skip
\begin{lstlisting}[style=coloredverbatimError]
fn foo (x: i32) {
    x = 1;
}

fn main () { foo (1); }
\end{lstlisting}

\begin{lstlisting}[style=bashVerb, breaklines=true, escapechar=~]
Error[E4156] : x is a value parameter, and cannot be used as a lvalue
 --> test_resources/diagnostics/test59.yr:(2,5)
 2  ~\lstbox{\textSFxi{}}~     x = 1;
    ~\lstbox{\textSFv{}}~     ^ ^
    ~\lstbox{\textSFii{}}~~\lstbox{\textSFx{}}~~\lstbox{\textSFx{}}~~\lstbox{\textSFx{}}~~\lstbox{\textSFx{}}~~\lstbox{\textSFx{}}~~\lstbox{\textSFx{}}~~\lstbox{\textSFx{}}~
\end{lstlisting}

\subsubsection*{How to fix it}

Declare the parameter \tokennolst{ref} if the caller's value has to be written, or bind the value to a local variable.

\errorcode{E4157}{value of type \errhole{} is not a lvalue}
\label{sec:codes:e4157}

Raised during \textbf{validation}, from \texttt{ValidateErrorMessage::\allowbreak{}NOT\_\allowbreak{}A\_\allowbreak{}LVALUE\_\allowbreak{}TYPE} (\textit{src/\allowbreak{}ymirc/\allowbreak{}semantic/\allowbreak{}validator/\allowbreak{}errors.yr}).

\subsubsection*{What it means}

A value of this type never has storage of its own -- a temporary, or the result of an expression.

\subsubsection*{Example}

\textit{test\_\allowbreak{}resources/\allowbreak{}function/\allowbreak{}test2.yr}:

%% check: skip
\begin{lstlisting}[style=coloredverbatimError]
fn foo (mut x : [mut i32]);

fn main () {
    let a = copy [1];
    foo (alias a);
}
\end{lstlisting}

\begin{lstlisting}[style=bashVerb, breaklines=true, escapechar=~]
Error[E4157] : value of type [i32] is not a lvalue
 --> test_resources/function/test2.yr:(5,10)
 5  ~\lstbox{\textSFxi{}}~     foo (alias a);
    ~\lstbox{\textSFv{}}~          ^^^^^ ^
    ~\lstbox{\textSFii{}}~~\lstbox{\textSFx{}}~~\lstbox{\textSFx{}}~~\lstbox{\textSFx{}}~~\lstbox{\textSFx{}}~~\lstbox{\textSFx{}}~~\lstbox{\textSFx{}}~~\lstbox{\textSFx{}}~
\end{lstlisting}

\subsubsection*{How to fix it}

Bind it to a variable and use that, cf. \hyperref[sec:codes:e4151]{E4151}.

\errorcode{E4165}{type \errhole{} is not a movable type}
\label{sec:codes:e4165}

Raised during \textbf{validation}, from \texttt{ValidateErrorMessage::\allowbreak{}NOT\_\allowbreak{}MOVABLE} (\textit{src/\allowbreak{}ymirc/\allowbreak{}semantic/\allowbreak{}validator/\allowbreak{}errors.yr}).

\subsubsection*{What it means}

\tokennolst{move} was applied to a type that is not movable. Moving is defined only for types that carry an identity worth transferring.

\subsubsection*{Example}

\textit{test\_\allowbreak{}resources/\allowbreak{}diagnostics/\allowbreak{}test60.yr}:

%% check: skip
\begin{lstlisting}[style=coloredverbatimError]
fn main () {
    let x = 1;
    let _ = move x;
}
\end{lstlisting}

\begin{lstlisting}[style=bashVerb, breaklines=true, escapechar=~]
Error[E4165] : type i32 is not a movable type
 --> test_resources/diagnostics/test60.yr:(3,13)
 3  ~\lstbox{\textSFxi{}}~     let _ = move x;
    ~\lstbox{\textSFv{}}~             ^^^^
    ~\lstbox{\textSFii{}}~~\lstbox{\textSFx{}}~~\lstbox{\textSFx{}}~~\lstbox{\textSFx{}}~~\lstbox{\textSFx{}}~~\lstbox{\textSFx{}}~~\lstbox{\textSFx{}}~~\lstbox{\textSFx{}}~
\end{lstlisting}

\subsubsection*{How to fix it}

Copy the value instead, or make the type movable.

\errorcode{E4167}{no copy exists for type \errhole{}}
\label{sec:codes:e4167}

Raised during \textbf{validation}, from \texttt{ValidateErrorMessage::\allowbreak{}NO\_\allowbreak{}COPY\_\allowbreak{}EXIST} (\textit{src/\allowbreak{}ymirc/\allowbreak{}semantic/\allowbreak{}validator/\allowbreak{}errors.yr}).

\subsubsection*{What it means}

\tokennolst{copy} was applied to a type that defines no copy.

\subsubsection*{Example}

\textit{test\_\allowbreak{}resources/\allowbreak{}diagnostics/\allowbreak{}test115.yr}:

%% check: skip
\begin{lstlisting}[style=coloredverbatimError]
// a copy of a value that holds nothing to copy
fn main () {
    let _ = copy 1;
}
\end{lstlisting}

\begin{lstlisting}[style=bashVerb, breaklines=true, escapechar=~]
Error[E4167] : no copy exists for type i32
 --> test_resources/diagnostics/test115.yr:(3,13)
 3  ~\lstbox{\textSFxi{}}~     let _ = copy 1;
    ~\lstbox{\textSFv{}}~             ^^^^
    ~\lstbox{\textSFxi{}}~
    = when validating test115::main -- (test_resources/diagnostics/test115.yr:(2,4))
    ~\lstbox{\textSFii{}}~~\lstbox{\textSFx{}}~~\lstbox{\textSFx{}}~~\lstbox{\textSFx{}}~~\lstbox{\textSFx{}}~~\lstbox{\textSFx{}}~~\lstbox{\textSFx{}}~~\lstbox{\textSFx{}}~
\end{lstlisting}

\subsubsection*{How to fix it}

Define the copy operation on the type, or pass the value without copying if a borrow is enough.

\errorcode{E4247}{aliasing the value of type \errhole{} to create a constant borrowing is prohibited}
\label{sec:codes:e4247}

Raised during \textbf{validation}, from \texttt{ValidateErrorMessage::\allowbreak{}UNECESSARY\_\allowbreak{}ALIAS} (\textit{src/\allowbreak{}ymirc/\allowbreak{}semantic/\allowbreak{}validator/\allowbreak{}errors.yr}).

\subsubsection*{What it means}

\tokennolst{alias} was used to produce a constant borrow. Aliasing exists to carry mutability across a borrow; producing a constant one from it says the borrow is not needed.

\subsubsection*{Example}

\textit{test\_\allowbreak{}resources/\allowbreak{}function/\allowbreak{}test3.yr}:

%% check: skip
\begin{lstlisting}[style=coloredverbatimError]
fn foo (x : [i32]);

fn main () {
    let dmut a = copy [1];
    foo (alias a);
}
\end{lstlisting}

\begin{lstlisting}[style=bashVerb, breaklines=true, escapechar=~]
Error[E4247] : aliasing the value of type [i32] to create a constant borrowing is prohibited
 --> test_resources/function/test3.yr:(5,10)
 5  ~\lstbox{\textSFxi{}}~     foo (alias a);
    ~\lstbox{\textSFv{}}~          ^^^^^  note: for parameter x: [i32]
    ~\lstbox{\textSFii{}}~~\lstbox{\textSFx{}}~~\lstbox{\textSFx{}}~~\lstbox{\textSFx{}}~~\lstbox{\textSFx{}}~~\lstbox{\textSFx{}}~~\lstbox{\textSFx{}}~~\lstbox{\textSFx{}}~
\end{lstlisting}

\subsubsection*{How to fix it}

Drop the \tokennolst{alias}: a constant borrow is what passing the value already gives.

\errorcode{E4248}{aliasing the value of type \errhole{} to call begin and end iterator constant methods is useless}
\label{sec:codes:e4248}

Raised during \textbf{validation}, from \texttt{ValidateErrorMessage::\allowbreak{}UNECESSARY\_\allowbreak{}ALIAS\_\allowbreak{}CPTR\_\allowbreak{}LOOP} (\textit{src/\allowbreak{}ymirc/\allowbreak{}semantic/\allowbreak{}validator/\allowbreak{}errors.yr}).

\subsubsection*{What it means}

\tokennolst{alias} was used to call the constant \tokennolst{begin} and \tokennolst{end} iterator methods of a loop, which do not need the mutability the alias carries.

\subsubsection*{How to fix it}

Drop the \tokennolst{alias} on the loop subject.

\errorcode{E4251}{the move has no effect}
\label{sec:codes:e4251}

Raised during \textbf{validation}, from \texttt{ValidateErrorMessage::\allowbreak{}UNECESSARY\_\allowbreak{}MOVE} (\textit{src/\allowbreak{}ymirc/\allowbreak{}semantic/\allowbreak{}validator/\allowbreak{}errors.yr}).

\subsubsection*{What it means}

\tokennolst{move} was applied where the source is not used afterwards anyway, so the move changes nothing.

\subsubsection*{Example}

\textit{test\_\allowbreak{}resources/\allowbreak{}regressions/\allowbreak{}move\_\allowbreak{}unnecessary\_\allowbreak{}ctor\_\allowbreak{}call/\allowbreak{}test1.yr}:

%% check: skip
\begin{lstlisting}[style=coloredverbatimError]
entity A {
    pub self () {}

    __dtor (mut self) {}
}

fn foo ()-> A {
    move A ()
}

fn main () {
    let _ = foo ();
}
\end{lstlisting}

\begin{lstlisting}[style=bashVerb, breaklines=true, escapechar=~]
Error[E4251] : the move has no effect
 --> test_resources/regressions/move_unnecessary_ctor_call/test1.yr:(8,5)
 8  ~\lstbox{\textSFxi{}}~     move A ()
    ~\lstbox{\textSFv{}}~     ^^^^   ^
    ~\lstbox{\textSFii{}}~~\lstbox{\textSFx{}}~~\lstbox{\textSFx{}}~~\lstbox{\textSFx{}}~~\lstbox{\textSFx{}}~~\lstbox{\textSFx{}}~~\lstbox{\textSFx{}}~~\lstbox{\textSFx{}}~
\end{lstlisting}

\subsubsection*{How to fix it}

Drop the \tokennolst{move}.

\errorcode{E4252}{referencing the value has no effect}
\label{sec:codes:e4252}

Raised during \textbf{validation}, from \texttt{ValidateErrorMessage::\allowbreak{}UNECESSARY\_\allowbreak{}REFERENCE} (\textit{src/\allowbreak{}ymirc/\allowbreak{}semantic/\allowbreak{}validator/\allowbreak{}errors.yr}).

\subsubsection*{What it means}

A reference was taken where the value is used as a value anyway.

\subsubsection*{Example}

\textit{test\_\allowbreak{}resources/\allowbreak{}for\_\allowbreak{}loops/\allowbreak{}arrays/\allowbreak{}test13.yr}:

%% check: skip
\begin{lstlisting}[style=coloredverbatimError]
fn main () {
    let dmut a = [1, 2, 3, 4];
    for x in ref a {
        x;
    }
}
\end{lstlisting}

\begin{lstlisting}[style=bashVerb, breaklines=true, escapechar=~]
Error[E4252] : referencing the value has no effect
 --> test_resources/for_loops/arrays/test13.yr:(3,14)
 3  ~\lstbox{\textSFxi{}}~     for x in ref a {
    ~\lstbox{\textSFv{}}~              ^^^
    ~\lstbox{\textSFii{}}~~\lstbox{\textSFx{}}~~\lstbox{\textSFx{}}~~\lstbox{\textSFx{}}~~\lstbox{\textSFx{}}~~\lstbox{\textSFx{}}~~\lstbox{\textSFx{}}~~\lstbox{\textSFx{}}~
\end{lstlisting}

\subsubsection*{How to fix it}

Drop the reference.
